\scanpage{29}
élément de $K$, on voit que le choix ainsi fait permet d'écrire
\[
 \alpha_\infty(u)=\frac{1}{2\sqrt{-1}\pi}\,v
\]
où $v$ est le générateur canonique de $H^1(\mathbf P^1,\Omega^1)$.

Pour les choix différents de topologie sur $K$ faisant de $K$ un corps valué complet arch., on trouve ainsi tous les éléments \unsure{primitifs|possibles} de $H^1(\mathbf P^1,\Omega^1)$.

Si $K$ est un corps \deleted{\unsure{top.}} isom à $\mathbf R$ [l'isom étant d'ailleurs \uline{unique}], alors on définit encore, par choix de $\mathbf C$ comme \deleted{\unsure{complexe}} clôture algébrique de $\mathbf R$, un foncteur $\xi=\xi_K$ à valeurs dans les modules sur $\mathbf Z$, \deleted{et \unsure{des isom}} \inserted{mais avec op. de $\mathbf Z/2\mathbf Z=\operatorname{Gal}(\overline{\mathbf R}/\mathbf R)$}. Le sous \unsure{$\mathbf Z$|$\mathbf Q$}-réseau \inserted{$\xi(M)$} que nous obtenons dans $\displaystyle\prod_\ell M_\ell\times M_{\infty\mathbf C}$ \inserted{\unsure{relatif :} $\mathbf R\hookrightarrow\mathbf C$} est alors stable par ce groupe $\mathbf Z/2\mathbf Z$

\mspage{28}
